%! TeX program = lualatex
\documentclass[../main.tex]{subfiles}
\begin{document} \section{Problem-solving strategies for the midterm exam}

Here are some exam writing strategies.

\begin{enumerate}
  \item Aim to correctly complete problems you can do \emph{as if there is no time pressure}.
  \begin{itemize}
    \item Avoid misreading the question. Highlight keywords in the problem statement.
    \item Avoid arithmetic and basic algebra mistakes.
  \end{itemize}

  \item Pay attention to details. Highlight easy-to-miss details. 

  \item Remember all techniques. The following are often forgotten.
    \begin{itemize}
      \item The limit definition of vertical asymptotes on page \pageref{def:vertical-asymptote}.

      \item The general change of variable formula for limits on page \pageref{sec:change-of-variable}.

      \item The substitution \(u = -x\) for limits at \(-\infty\) and limits at infinity with exponential functions.

      \item Recognizing that a limit is the derivative of a function (Problem 7 in Quiz 3).

      \item The technique to find the derivative of an inverse function on page \pageref{technique:implicit-differentiation-inverse}.
    \end{itemize}
    
  \item It is okay to learn problem-solving with others. However, you also need to practice solving problems on your own.
\end{enumerate}

\bigskip

\faStar{} Lastly, hard problems can be solved by \hlattn{thinking systematically} to break up the problem into smaller, more manageable parts.  Other than the systematic applications of \emph{a group of related techniques}, here are some big picture questions to ask yourself when solving a problem.
\begin{itemize}
  \item What is the form of the final answer?
  \item What missing functions or quantities do I need to calculate?
  \item How can I calculate the missing functions or quantities? Can I use a routine technique or do I need to resort to the definitions?
  \item Have I seen a similar problem before? This requires you to actively recall the examples, exercises and quiz problems you have done before. 
\end{itemize}
\clearpage

All problems below are similar to past exam problems with less than \(60\%\) success rate.

\begin{example}
  Let \(f(x) = \begin{cases} x^{2} \cos\left(\frac{1}{x^{2}} \right) &\text{if } x < 0\\ 2x &\text{if } x \geq 0 \end{cases}\). Which statement is true?

  \begin{enumerate}[label=\Alph*.]
    \item \(f\) is not defined at \(0\).
    \item \(f\) is not continuous at \(0\).
    \item \(f\) is continuous at \(0\) but not differentiable at \(0\).
    \item \(f\) is differentiable at \(0\) but not continuous at \(0\).
    \item \(f\) is differentiable at \(0\).
  \end{enumerate} 
  
  \blanks*
\end{example}


\begin{example}
  Evaluate \(\lim_{x \to 0^{-}} e^{\frac{1}{\sin(x)}}\).

  \blanks*
\end{example}
\clearpage

\begin{exercise}
  Find \(\frac{dy}{dx}\) given that \(y^{2} + 3x^{2}\sqrt{y} = 5\).
  
  \blanks*
\end{exercise}

\begin{exercise}
  Let \(f(x) = \frac{1}{|x-1|} + |x^{3}-x^{2}|\). Find \(f'(0)\).

  \blanks*
\end{exercise}

\begin{example}
  Evaluate \(\lim_{h \to 0} \frac{3^{h+1} - 3}{h}\).
  
  \blanks*
\end{example}
\clearpage

\begin{exercise}
  Find \(x\) at which \(x^{2} - xy + y^{2} = 3\) has a vertical tangent line.
  
  \blanks*
\end{exercise}

\begin{exercise}
  Find \(x\) at which \(\arctan(x)\) has a vertical tangent line.

  Legit or fishy? Let \(y = \arctan(x)\). The question translates to solving for \(x\) given \(\frac{dx}{dy} = 0\).
  \[
    \frac{d}{dy} y = \frac{d}{dy} \arctan(x)
    \quad\longrightarrow\quad
    1 = \frac{1}{1 + x^{2}}
  \]
  The equation is true at \(x = 0\). So \(\arctan(x)\) has a vertical tangent line at \(x = 0\).
\end{exercise}
\clearpage

\begin{exercise} \label{ex:tangent-parallel-to-line}
  Find all points \((x,y)\) at which the tangent line to the curve defined by the equation
  \[
    x^{2} + xy + y^{2} = 1
  \]
  is parallel to the line \(y = -2x + 3\). 
  
  \blanks*
\end{exercise}
\clearpage

\begin{exercise}
  Evaluate \(\frac{d}{dx} \sin(\sec^{-1}(\sqrt{\pi}))\).
  
  \blanks*
\end{exercise}

\begin{exercise}
  Let \(f(x) = \cos(3x) + (x+1)^{2}\). Find \(f^{(5)}(0)\). 
  
  \blanks*
\end{exercise}

\begin{exercise}
  Suppose \(f(x) = x + 2^{x}\). Assume \(f^{-1}(x)\) exists. Find \(\left(f^{-1}\right)'(6)\).
  
  \blanks*
\end{exercise}

\begin{exercise}
  Let \(L\) be a line tangent to the curve \(y = x^{3} + x\). If \(L\) is parallel to \(y = x\), which of the following could be an equation of \(L\)?

  \begin{enumerate}[label=\Alph*.]
    \item \(y = x + 1\)
    \item \(y = x - 1\)
    \item \(y = x\)
    \item \(y = 0\)
    \item \(y = 1\)
  \end{enumerate}

  \blanks{5}
\end{exercise}
\clearpage


\begin{exercise}
  The equation \(\left(\frac{x}{2}\right)^{2} + \left( \frac{y}{3} \right)^{2} = 1\) defines an ellipse.  In the first quadrant where \(x \ge 0\) and \(y \ge 0\), there is exactly one point \((x_{0}, y_{0})\) on the ellipse where the tangent line to the ellipse has slope \(m = -9/2\).  Use implicit differentiation to find this point.

  \blanks*
\end{exercise}
\clearpage

\end{document}